\subsection{Modules of finite presentation}

A presentation (or presentation of length 1) of an A-module E is an exact sequence

\[
\mathrm{L} _ {1} \rightarrow \mathrm{L} _ {0} \rightarrow \mathrm{E} \rightarrow 0\tag{13}
\]

of A-modules where \(\mathbf{L}_0\) and \(\mathbf{L}_1\) are free.

Every A-module E admits a presentation. Indeed, we know (II, p.~27, prop. 20) that there exists a surjective homomorphism \(u: \mathbf{L}_0 \to \mathbf{E}\) , where \(\mathbf{L}_0\) is free; if R is the kernel of \(u\) , there likewise exists a surjective homomorphism \(v: \mathbf{L}_1 \to \mathbf{R}\) where \(\mathbf{L}_1\) is free. If one regards \(v\) as a homomorphism of \(\mathbf{L}_1\) to \(\mathbf{L}_0\) , the sequence \(\mathbf{L}_1 \xrightarrow{v} \mathbf{L}_0 \xrightarrow{u} \mathbf{E} \longrightarrow 0\) is exact by definition, whence our assertion.

If \(\rho: A \to B\) is a ring homomorphism, every presentation (13) of \(E\) furnishes a presentation of \(E_{(B)} = B \otimes_A E\) :

\[
\mathbf {B} \otimes_ {\mathbf {A}} \mathrm{L} _ {1} \rightarrow \mathbf {B} \otimes_ {\mathbf {A}} \mathrm{L} _ {0} \rightarrow \mathbf {B} \otimes_ {\mathbf {A}} \mathrm{E} \rightarrow 0\tag{14}
\]

by virtue of II, p.~58, prop. 5 and of the fact that \(\mathbf{B} \otimes_{\mathbf{A}} \mathbf{L}\) is a free B-module when L is free.

One says that a presentation (13) of a module E is finite if the free modules \(L_{0}\) and \(L_{1}\) have finite bases. It is clear that if the presentation (13) is finite, so is the presentation (14). One says that E is an A-module of finite presentation if it admits a finite presentation.

PROPOSITION 5. --- (i) Every module admitting a finite presentation is finitely generated.

\begin{enumerate}
\def\labelenumi{(\roman{enumi})}
\setcounter{enumi}{1}
\item
  If A is a left Noetherian ring, every finitely generated A-module admits a finite presentation.
\item
  Every finitely generated projective module admits a finite presentation.
\end{enumerate}

Assertion (i) follows trivially from the definitions. Suppose A is left Noetherian and E is finitely generated. There then exists a surjective homomorphism \(u: L_0 \to E\) , where \(L_0\) is a free A-module having a finite basis; the kernel R of \(u\) is finitely generated, hence there is a surjective homomorphism \(v: L_1 \to R\) where \(L_1\) is free with finite basis, and the exact sequence \(L_1 \xrightarrow{v} L_0 \xrightarrow{u} E \longrightarrow 0\) is a finite presentation of E; whence (ii).

Finally, suppose that E is a finitely generated projective module; it is then a direct summand of a finitely generated free module \(L_{0}\) (II, p.~40, cor. 1); the kernel R of the surjective homomorphism \(L_{0} \rightarrow E\) is then isomorphic to a quotient of \(L_{0}\) , hence is finitely generated, and one finishes as above.

PROPOSITION 6. --- Let A be a ring, E an A-module of finite presentation. For every exact sequence

\[
0 \longrightarrow \mathrm{F} \stackrel {j} {\longrightarrow} \mathrm{G} \stackrel {p} {\longrightarrow} \mathrm{E} \longrightarrow 0
\]

where G is finitely generated, the module F is finitely generated.

Let \(L_{1} \xrightarrow{r} L_{0} \xrightarrow{s} E \longrightarrow 0\) be a finite presentation; if \((e_{i})\) is a basis of \(L_{0}\) , there exists for each i an element \(g_{i} \in G\) such that \(p(g_{i}) = s(e_{i})\) ; the homomorphism \(u : L_{0} \to G\) such that \(u(e_{i}) = g_{i}\) for all i is therefore such that \(s = p \circ u\) . Since \(s \circ r = 0\) , we have \(u(r(L_{1})) \subset \operatorname{Ker} p\) , and since Ker p is isomorphic to F, we see that there is a homomorphism \(v : L_{1} \to F\) such that the diagram

\[
\begin{array}{c c c c c c c} L _ {1} & \xrightarrow {r} & L _ {0} & \xrightarrow {s} & E & \longrightarrow & 0 \\ v \Big \downarrow & & u \Big \downarrow & & ^ {1 _ {E}} \Big \downarrow \\ F & \xrightarrow {} _ {j} & G & \xrightarrow {} _ {p} & E & \longrightarrow & 0 \end{array}
\]

is commutative. Since j is injective and s surjective, one can apply Proposition 2 of X, p.~4, in other words there is an exact sequence

\[
\operatorname{Ker} 1 _ {\mathrm{E}} \xrightarrow {d} \text { Coker } v \longrightarrow \text { Coker } u \longrightarrow \text { Coker } 1 _ {\mathrm{E}}.
\]

This shows that Coker v is isomorphic to \(\mathrm{G}/u(\mathrm{L}_{0})\) , which is finitely generated by hypothesis. Moreover, we have the exact sequence

\[
0 \rightarrow v (\mathbf {L} _ {1}) \rightarrow \mathrm{F} \rightarrow \text { Coker } v \rightarrow 0
\]

and since \(v(\mathbf{L}_1)\) and Coker \(v\) are finitely generated, so is F (II, p.~17, cor. 5).

PROPOSITION 7. --- Let M be an A-module. There exists a right-filtered ordered set I and an inductive system of finitely presented A-modules (M \(_{\alpha}\) , φ \(_{\beta\alpha}\) ) relative to I such that M is isomorphic to lim M \(_{\alpha}\) . If M possesses a generating family of n elements, one may assume that the same is true of the M \(_{\alpha}\) .

Consider a presentation

\[
\mathrm{A} _ {s} ^ {(\mathrm{K})} \xrightarrow {u} \mathrm{A} _ {s} ^ {(\mathrm{L})} \xrightarrow {v} \mathrm{M} \longrightarrow 0;
\]

let I be the set of pairs \(\alpha = (\mathbf{K}', \mathbf{L}')\) , where \(\mathbf{K}'\) (resp. \(\mathbf{L}'\) ) is a finite subset of \(\mathbf{K}\) (resp. L), such that \(u\) induces a map \(u_{\alpha}\) of the submodule \(\mathbf{A}_s^{\mathbf{K}'}\) of \(\mathbf{A}_s^{(\mathbf{K})}\) in the submodule \(\mathbf{A}_s^{\mathbf{L}'}\) of \(\mathbf{A}_s^{(\mathbf{L})}\) ; for \(\alpha \in \mathbf{I}\) , let \(\mathbf{M}_{\alpha}\) be the cokernel of \(u_{\alpha}\) and \(v_{\alpha}: \mathbf{A}_s^{\mathbf{L}'} \to \mathbf{M}_{\alpha}\) the canonical map, so that we have a commutative diagram with exact rows:

\[
\mathrm{A} _ {s} ^ {(\mathrm{K})} \xrightarrow {u} \mathrm{A} _ {s} ^ {(\mathrm{L})} \xrightarrow {i} \mathrm{M} \longrightarrow 0
\]

\[
i _ {\alpha} \uparrow , \quad j _ {\alpha} \uparrow , \quad f _ {\alpha} \uparrow
\]

\[
\mathrm{A} _ {s} ^ {\mathrm{K} ^ {\prime}} \xrightarrow {u _ {\alpha}} \mathrm{A} _ {s} ^ {\mathrm{L} ^ {\prime}} \xrightarrow {v _ {\alpha}} \mathrm{M} _ {\alpha} \longrightarrow 0,
\]

where \(i_{\alpha}\) and \(j_{\alpha}\) are the canonical injections, and where \(f_{\alpha}\) is deduced from \(j_{\alpha}\) by passing to quotients. Order the set I by the relation

\[
\alpha = \left(\mathrm{K} ^ {\prime}, \mathrm{L} ^ {\prime}\right) \leqslant \beta = \left(\mathrm{K} ^ {\prime \prime}, \mathrm{L} ^ {\prime \prime}\right) \quad \text { if } \quad \mathrm{K} ^ {\prime} \subset \mathrm{K} ^ {\prime \prime}, \quad \mathrm{L} ^ {\prime} \subset \mathrm{L} ^ {\prime \prime};
\]

for \(\alpha \leqslant \beta\) , let \(\varphi_{\beta \alpha}: \mathbf{M}_{\alpha} \to \mathbf{M}_{\beta}\) be the homomorphism deduced by passing to quotients from the inclusion of \(\mathbf{A}_s^{L'}\) to \(\mathbf{A}_s^{L''}\) . One then immediately verifies that the ordered set I is filtered, and that \((\mathbf{M}_{\alpha}, \varphi_{\beta \alpha})\) is an inductive system of A-modules and that \((\varphi_{\alpha})\) is an inductive system of A-homomorphisms. Passing to the inductive limit, one obtains a commutative diagram

\[
\begin{array}{c c c c c} \mathbf {A} _ {s} ^ {(\mathbf {K})} & \stackrel {{u}} {{\longrightarrow}} & \mathbf {A} _ {s} ^ {(L)} & \stackrel {{v}} {{\longrightarrow}} & \mathbf {M} \\ i \uparrow & & j \uparrow & & \varphi \uparrow \\ \underline {{\lim}} \mathbf {A} _ {s} ^ {\mathbf {K} ^ {\prime}} & \longrightarrow & \underline {{\lim}} \mathbf {A} _ {s} ^ {\mathbf {L} ^ {\prime}} & \longrightarrow & \underline {{\lim}} \mathbf {M} _ {\alpha} \longrightarrow 0; \end{array}\tag{15}
\]

The rows of this diagram are exact (II, p.~91, prop. 3); since i and j are bijective, φ is also bijective (X, p.~7, cor. 3), whence the proposition.

